\section{Interpreting calibration across settings}
\subsection{Which comparison can establish improvement?}
A calibration rule is evaluated against a stated financial preference. In this study, improvement in alignment means a smaller $d_t^r$ or $g_t^r$ on the same estimated state. It does not mean that the calibrated strategy earns a higher return in the next month. The realised ranking can reverse because the target's estimated mean and covariance are imperfect forecasts. This distinction is especially relevant when the target itself uses a noisy expected-return estimate and a competing policy happens to reduce exposure to that estimate.

The primary distances put the scale of the comparison in context. Raw second moments give mean Euclidean distance 0.00843; unconstrained scalar calibration reduces it to 0.00540; reference-based risk-only calibration gives 0.00510. Joint calibration and the MVPI adaptation agree with the target to numerical precision. The log-minus-variance adaptation has distance 0.01286. These are distances between ten-dimensional weight vectors, not percentages of portfolio wealth lost. Fitted regret supplies the corresponding objective scale: the raw reward gap is only 0.0148 monthly basis points. A statistically visible allocation difference can therefore be economically small under the fitted preference.

The historical MV score answers a different question. Joint calibration is 1.315 monthly basis points below raw in the primary sample, while its interval spans both signs. That result does not support a claim that the calibrated policy is a superior investment strategy. Nor does it disprove the calibration identity: subsequent realised returns are absent from that identity. Minimum variance's better realised MV and Sharpe suggest that a strong practical implementation would need to address estimation quality alongside reward design.

\subsection{What the robustness grid tests}
The first block crosses three risk coefficients, $\varphi\in\{2,5,10\}$, with three turnover charges, $\lambda\in\{0,0.001,0.0025\}$, at the primary estimation window and shrinkage. The second crosses $L\in\{60,120,240\}$ with $h\in\{0,0.1,0.5\}$ at the primary risk and cost coefficients. Their shared primary setting is counted once, giving 17 distinct specifications. This is a union of two grids rather than the full cross-product of every choice. It probes economically relevant perturbations without pretending to span all market models.

The maximum setting-mean distance for joint calibration is $1.54\times10^{-7}$, compared with raw mean distances ranging from approximately 0.00157 to 0.03682. This supports the identity across the evaluated covariance estimates, constraints and cost levels. It does not supply 17 independent statistical replications: all settings reuse the same asset universe and overlapping return observations. The robustness runs compare the target, raw, scalar, risk-only, joint, equal-weight and minimum-variance policies. The MVPI and logarithmic comparators are evaluated in the primary study, rather than silently treated as rerun throughout this grid.

At zero transaction cost the risk-only and joint reward definitions coincide. With positive cost, the difference depends on the target expected return through $A=1+2\varphi y^*$. Constraints can make distinct objectives select the same boundary allocation, so exact agreement at an individual state need not establish equivalence of the rewards at all states. Conversely, a small discrepancy at the primary setting cannot be assumed to remain small when risk, cost or estimation changes.

\subsection{Optimisation effort and reward fidelity}
The finite-budget comparison assigns independent logits to each historical state. Softmax enforces strictly positive weights at finite logits, so exact boundary allocations can be approached without being attained. This helps explain why a high-accuracy constrained solver and a finite-budget softmax optimiser need not have equal residual errors even on the same objective. All three trained objectives use identical initialisation distributions, seeds and step budgets; the diagnostic compares numerical effort under that common parameterisation.

The own-reward gap measures whether the optimiser solved the reward it was given. The signed transfer term in the target decomposition instead measures how the approximate policy changes the target value relative to the exact reward optimum. It can have either sign. Adding own-reward suboptimality directly to target-unit design regret would mix different objective scales and would not reproduce the total target gap. At 5,000 steps, mean own-reward gaps remain around $4\times10^{-6}$, larger than the exact raw design gap of $1.48\times10^{-6}$. Improving the optimiser can therefore be the larger immediate numerical intervention, while calibration addresses the distinct discrepancy that persists at exact optimisation.
